%% Carta de presentacion para EAAI. Lleva identidad: no es el
%% manuscrito anonimo. La firma el autor de correspondencia.
\documentclass[11pt]{article}
\usepackage[margin=2.2cm]{geometry}
\usepackage{lmodern}
\usepackage[T1]{fontenc}
\usepackage[utf8]{inputenc}
\usepackage{parskip}
\setlength{\parskip}{0.55em}
\pagestyle{empty}

\begin{document}

\noindent
Hern\'an D\'iaz Rodr\'iguez\\
Department of Computing, University of Oviedo\\
Gij\'on, Spain\\
diazhernan@uniovi.es

\vspace{1em}
\noindent 1 September 2026

\vspace{1em}
\noindent To the Editors,\\
\emph{Engineering Applications of Artificial Intelligence}

\vspace{1em}

We submit for your consideration the manuscript \emph{Deep
Reinforcement Learning for the Interval Job Shop Scheduling Problem:
A Comparison with Genetic Programming Hyper-Heuristics across
Inference Budgets}, as a research paper.

The interval job shop scheduling problem models each uncertain
processing time as an interval between known bounds, and the
literature solves it with metaheuristics that search anew on every
instance. The contribution in artificial intelligence is twofold. The
first is a constructive deep reinforcement learning policy for the
problem whose features are dimensionless and whose encoder is
permutation invariant, so that a single network trained on four
$20\times15$ instances schedules any other size without retraining or
architectural change; it improves on every hand-crafted dispatching
rule on all seventy instances of the benchmark, six of whose seven
size classes it never saw. The second is a controlled comparison
between the two learned constructive paradigms available for the
problem, thirty trained artifacts of each family evaluated in one
harness, under one selection criterion and at matched inference
budgets. That comparison finds no dominance: which learner leads is
decided by the inference budget and, on instances neither family has
seen, by the instance size class. The engineering application is
production scheduling under uncertain processing times, and the paper
reports the training and inference costs that a deployment decision
would need.

The work suits this journal because its result is an engineering
characterization rather than a leaderboard entry: it says when a
compact, readable evolved rule is preferable and when a policy's
sampling distribution earns its higher cost.

We confirm that the manuscript is original, that it has not been
published elsewhere, and that it is not under consideration by another
journal. We disclose one related work: a companion study by the same
first author on genetic programming hyper-heuristics for this problem
is under review elsewhere and available as a preprint (SSRN,
doi:10.2139/ssrn.7214286). The present manuscript is distinct in both
contributions above, and that study enters it only as the published
source of the thirty evolved rules it compares against.

The instances, source code, experiment records and trained
checkpoints supporting the study are openly available in Zenodo at
doi:10.5281/zenodo.21970431. A statement disclosing the use of
generative AI in the preparation of the manuscript appears directly
before the references, as the journal's policy requires. In keeping
with the double-anonymous review process, author identity appears
only in the title page and in this letter.

Thank you for considering our work.

\vspace{0.7em}
\noindent Yours sincerely,

\vspace{1.2em}
\noindent Hern\'an D\'iaz Rodr\'iguez\\
Corresponding author, on behalf of both authors

\end{document}
